\section{Discussion}
\label{sec:mech}

Coherence is most useful when the mapping task requires separating surfaces with similar backscatter but different temporal stability. Bright returns from buildings and forest can overlap in intensity, while persistent structures and changing canopy scatterers generate distinct repeat-pass coherence distributions~\cite{Sica2019,Aghababaei2026}. This distinction gives the observable a clear role in mixed rural landscapes, where small settlements are embedded in vegetation. The class distributions and error maps connect that scattering interpretation to the recovery of built-up land cover. For deployment, the relevant question is therefore whether temporal stability separates the classes that matter to the application. A landscape containing many rigid structures next to forest presents a different classification problem from one dominated by agricultural classes with similar 12-day decorrelation.

Recent interferometric studies help explain why the coherence contrast is informative. Forest scattering models link repeat-pass coherence to canopy and ground contributions~\cite{Aghababaei2026}. Crop growth and soil-moisture changes alter interferometric observations through vegetation development and dielectric properties~\cite{Yuan2024,Ao2025}, while random motion produces temporal decorrelation~\cite{Chen2026}. The resulting sensitivity to surface change also underlies the amplitude--coherence inputs used by DANI-NET for change detection~\cite{Costa2025}. Here, that temporal sensitivity supplies a class-discriminating feature at a single repeat interval.

The evaluation objective determines how that class separation translates into practical value. Sample overall accuracy weights each class-recall change by its prevalence in the evaluation sample. Balanced accuracy assigns every class equal weight. A coherence feature can consequently contribute strongly to settlement detection while producing a smaller change in scene-wide accuracy. For an application that must retain small settlements, built-up recall and its spatial support are directly relevant alongside balanced accuracy. For land-cover inventories, sample overall accuracy describes the aggregate outcome on the evaluated support. Reporting both connects the same classifier to these different objectives. The spring and autumn evaluations further motivate tracking class-level performance across acquisition windows: the mapping benefit depends on the distribution of errors among classes as well as the separability of built-up land cover.

At a two-acquisition budget, coherence extracts temporal information from the complex image pair already used by the intensity representation. The compact $I_8{+}g_{12}$ representation offers a practical starting point for deployment, retaining the observed discrimination with few feature channels. The ablation supports using the raw scalar directly; spatial coherence derivatives add processing steps without improving the evaluated classifiers. Acquiring and processing a suitable interferometric pair is thus the central operational requirement. Once coherence is available, deployment can begin with the compact representation and compare its performance with the richer intensity features on the intended source--target track pair. Pair-specific validation remains useful because acquisition configuration determines the classifier's operating accuracy. The broader methodological contribution is a comparison in which additional feature information is evaluated at a fixed observation budget. Applying this same design to other budgets would make acquisition planning and representation choice part of a common quantitative assessment.

Coherence estimation provides a concrete avenue for extending this representation. Recent methods select similarly decorrelated neighborhoods~\cite{Yao2024}, reduce low-coherence estimation bias~\cite{LiangH2026}, and estimate interferometric parameters with self-supervised learning~\cite{Geara2025}. MTCSAR uses a coherence prior to guide multitemporal SAR denoising~\cite{LiuJ2025}. The observation model also matters: burst-mode phase separation~\cite{Pulella2024} and analyses of range misregistration~\cite{Yang2026} connect interferometric measurements to acquisition and processing geometry. An extension can compare these estimators at the same acquisition budget and measure the resulting change in class discrimination.

Acquisition planning can build on coherence proxies that describe baseline and seasonal variation~\cite{Xu2025}, while scatterer-specific stochastic models address time-varying interferometric observation quality~\cite{Brouwer2025}. For land-cover products, field-boundary mapping offers an object-level complement to pixel classification~\cite{Celik2026}. Spatial and temporal validation of vegetation estimates likewise connects representation choice to deployment conditions~\cite{WangX2025}. Together, these developments suggest a practical sequence: select the repeat-pass pair, estimate the additional observable, and evaluate its contribution on the spatial units relevant to the mapping product.

The present evidence concerns orbital and seasonal deployment within one region. Acquisition geometry and date vary together, so the class-separation result describes their combined operating conditions. The built-up sample contains 165 pixels from 16 objects, and the reference maps predate the SAR acquisitions; support and label-sensitivity diagnostics are provided in the supplement. Geographic replication with contemporaneous labels is the next test of the measured effect sizes.

\section{Conclusion}

A single 12-day interferometric coherence measurement adds substantial class-discriminating information to two-acquisition Sentinel-1 land-cover classification. Against a spatially tuned 28-D intensity baseline, it increased cross-track balanced accuracy by $6.30\pp$ for RF and $6.99\pp$ for SVC, with overall-accuracy gains of $1.95\pp$ and $1.64\pp$. The four-track comparison ties this improvement to stronger built-up/forest separation. Comparable in-track gains identify temporal stability as useful class information across deployment configurations. The compact nine-dimensional representation retains nearly all of the performance achieved by the larger augmented feature set, providing a simple implementation of this finding.

The study quantifies coherence's marginal value while holding the acquisition budget, training support, and model-selection procedure constant. Its representation ladder links feature choice to class-level error recovery, and the paired spatial evaluation measures the gain across ordered orbital transfers. Autumn replication retained balanced-accuracy gains of $5.03\pp$ and $5.23\pp$ for the two learners. Together, these experiments provide an empirical basis for including raw coherence in low-observation land-cover mapping and a reproducible framework for evaluating additional SAR observables against competitive intensity representations.

